\Needspace{8\baselineskip}\originalchapter{2002 年波动与动力学作业}
\originalsection{2002 PS5：极化电流、群速与离子声}
\begin{exercise}\label{ex:hist02-5-1}
PS5 原题1（有官方答案）。从离子极化漂移求横向线性“电导” $\sigma_\perp$，用它重新推出平行传播 Alfv\'en 波；解释为什么理想 MHD 的 $\mathbf E+\mathbf v\times\mathbf B=0$ 看似只给 $E\times B$ 漂移，却仍包含极化效应。题卷 Gaussian 式为 $\sigma=-\ii\omega\rho c^2/B_0^2$。
\end{exercise}
\begin{solution}
编者独立解答。采用 $e^{\ii(kz-\omega t)}$ 和 SI，$\mathbf v_{pi}=m_i\partial_t\mathbf E_\perp/(eB_0^2)$，$\mathbf J_\perp=en_i\mathbf v_{pi}=\rho\partial_t\mathbf E_\perp/B_0^2$，故 $\sigma_\perp=-\ii\omega\rho/B_0^2$。Faraday 与 Amp\`ere 的横向分量（低频忽略位移电流）给 $k^2E_\perp=\ii\omega\mu_0J_\perp$，即 $\omega^2=k^2B_0^2/(\mu_0\rho)=k^2v_A^2$。Gaussian 的电流与场单位使式子多 $c^2$。理想 MHD 的 Ohm 式只是流体感应场约束，动量方程还含 $\rho\partial_t\mathbf v$；把 $\mathbf v=\mathbf E\times\mathbf B/B^2$ 代进动量方程，该惯性恰等价上述极化电流。若同时删去惯性，才会错误地消掉波。
\end{solution}

\begin{exercise}\label{ex:hist02-5-2}
PS5 原题2（有官方答案）。证明高斯包络 $A(x,0)=\Re[e^{\ii k_0x}e^{-x^2/(2\sigma^2)}]$ 有高斯平面波谱；对 $k_0\gg1/\sigma$ 且谱窄的色散关系 $\omega(k)$，证明包络以 $v_g=\omega'(k_0)$ 移动，讨论 $\omega''(k_0)$ 对展宽的作用，及真空光波的特例。
\end{exercise}
\begin{solution}
编者独立解答。完成平方的高斯积分给
\begin{equation}A(x,t)=\Re\left[\frac\sigma{\sqrt{2\pi}}\int_{-\infty}^\infty\exp\left(-\frac{\sigma^2(k-k_0)^2}{2}+\ii kx-\ii\omega(k)t\right)\dd k\right].\end{equation}
$t=0$ 时积分正好是题给初形。令 $q=k-k_0$、$\omega=\omega_0+v_gq+\frac12\omega''q^2+\cdots$，保留一阶时 \begin{equation}A=\Re\!\left[e^{\ii(k_0x-\omega_0t)}e^{-(x-v_gt)^2/(2\sigma^2)}\right].\end{equation}，快振荡相速和慢包络群速不同。保留二阶则复宽度为 $\sigma^2+\ii\omega''t$，真实包络标准差增至 $\sigma\sqrt{1+(\omega''t/\sigma^2)^2}$；谱内 $\omega''=0$ 才无此色散展宽。真空固定传播方向 $\omega=ck$ 在 $k>0$ 正满足；跨过 $k=0$ 时不能把 $\omega=c|k|$ 当全球线性。
\end{solution}

\begin{exercise}\label{ex:hist02-5-3}
PS5 原题3（有官方答案）。在无磁场、低频、准中性等离子体中，电子取 Boltzmann 绝热密度响应 $n_{e1}/n_0=e\phi/T_e$，冷离子按电场自由运动；推导 $\omega^2=k^2c_{IA}^2$，其中 $c_{IA}=\sqrt{T_e/m_i}$，并说出适用尺度。
\end{exercise}
\begin{solution}
编者独立解答。$-\ii\omega n_{i1}+\ii kn_0u_i=0$、$-\ii\omega m_iu_i=-\ii ke\phi$，所以 $n_{i1}/n_0=k^2e\phi/(m_i\omega^2)$。准中性 $n_{i1}=n_{e1}$ 给 $\omega^2=k^2T_e/m_i$。需要 $k\lambda_{De}\ll1$、电子足够快可建 Boltzmann 响应、离子热压 $T_i\ll T_e$；若保留 Poisson 但仍冷离子，$\omega^2=k^2c_{IA}^2/(1+k^2\lambda_{De}^2)$，说明短波偏离。
\end{solution}

\originalsection{2002 PS6：Fourier 核、熵与 Vlasov 极限}
\begin{exercise}\label{ex:hist02-6-1}
PS6 原题1（有官方答案）。取 Fourier 变换 $\widetilde f(\mathbf k)=\int e^{-\ii\mathbf k\cdot\mathbf x}f(\mathbf x)\dd^3x$，用积分及 Poisson 方程两种办法证明 $\widetilde{(1/r)}=4\pi/k^2$，并给 $\widetilde{\partial_{x_j}(1/r)}$。
\end{exercise}
\begin{solution}
编者独立解答。角积分后加收敛因子 $e^{-\epsilon r}$，$\int_0^\infty4\pi r e^{-\epsilon r}\sin(kr)/(kr)\dd r=4\pi/(k^2+\epsilon^2)$，取 $\epsilon\downarrow0$ 得分布意义的 $4\pi/k^2$。另一法是 $\nabla^2(1/r)=-4\pi\delta^3(\mathbf r)$，变换后 $-k^2\widetilde{(1/r)}=-4\pi$；边界要求排除任意 $\delta(\mathbf k)$ 齐次项。分部积分给 $\widetilde{\partial_{x_j}(1/r)}=\ii k_j4\pi/k^2$。题卷的分量写作 $\partial_x$，这里写任意 $j$ 以明其向量结构。
\end{solution}

\begin{exercise}\label{ex:hist02-6-2}
PS6 原题2（有官方答案）。计算 $f(\mathbf x)=(2\pi)^{-3}\int [4\pi/(k^2+k_0^2)]e^{-\ii\mathbf k\cdot\mathbf x}\dd^3k$，$k_0>0$。
\end{exercise}
\begin{solution}
编者独立解答。角积分化成 $f(r)=2/(\pi r)\int_0^\infty k\sin(kr)/(k^2+k_0^2)\dd k$，闭合上半平面或使用正弦积分 $\int_0^\infty k\sin(kr)/(k^2+k_0^2)\dd k=(\pi/2)e^{-k_0r}$，故 $f(r)=e^{-k_0r}/r$。另一核对是 $(\nabla^2-k_0^2)f=-4\pi\delta^3$；$r=0$ 处函数奇异，变换按分布解释。
\end{solution}

\begin{exercise}\label{ex:hist02-6-3}
PS6 原题3（有官方答案）。均匀系统在固定 $N=\int f\dd^3v$、$E=\int mv^2f/2\dd^3v$ 下极大化 $S=-\int f\ln f\dd^3v$，用变分推 Maxwell 分布，并说明为什么是极大值。
\end{exercise}
\begin{solution}
编者独立解答。$\delta(S+\alpha N+\beta E)/\delta f=-(\ln f+1)+\alpha+\beta mv^2/2=0$，故 $f=Ae^{-mv^2/(2T)}$，其中 $T=-1/\beta>0$，$A=N[m/(2\pi T)]^{3/2}$，$E=3NT/2$。在满足线性约束的变化上，$\delta^2S=-\int(\delta f)^2/f\dd^3v<0$，因此确为严格极大。若总动量也固定，需另一矢量乘子，结果是漂移 Maxwell；原题仅指定粒数与能量，故中心为零。
\end{solution}

\begin{exercise}\label{ex:hist02-6-4}
PS6 原题4（有官方答案）。把每个宏观粒子分成 $N$ 份，$q\mapsto q/N,m\mapsto m/N,n\mapsto Nn$。证明 Vlasov 方程及自洽场保持不变，但离散参数 $1/(n\lambda_D^3)$ 随 $N\to\infty$ 趋零。
\end{exercise}
\begin{solution}
编者独立解答。以归一化速度分布 $F=f/n$ 写 Vlasov：$\partial_tF+\mathbf v\cdot\nabla F+(q/m)(\mathbf E+\mathbf v\times\mathbf B)\cdot\nabla_vF=0$；$q/m$ 不变。电荷与质量密度分别 $(Nn)(q/N)=nq$、$(Nn)(m/N)=nm$，故 Maxwell 自洽源及惯性不变。电子 Debye 长度 $\lambda_D^2=\epsilon_0T/(nq^2)$ 在替换后变为 $N\lambda_D^2$，若把每个碎片热能温度保持同一速度方差，则粒子质量缩小使 $T\mapsto T/N$，两者正好抵消，$\lambda_D$ 不变。于是 $1/(n'\lambda_D'^3)=1/(Nn\lambda_D^3)\to0$；这指平均场颗粒噪声消失，不证明任何有限 $N$ 真实等离子体严格无碰撞。必须连温度的物理定义一起变换，不能在 $q,m,n$ 改变时暗留 $T$ 不变。
\end{solution}

\begin{exercise}\label{ex:hist02-6-5}
PS6 原题5（有官方答案）。课程 Laplace 约定及其逆变换为\begin{equation}g_\omega=\int_0^\infty e^{\ii\omega t}g(t)\dd t,\qquad g(t)=\frac1{2\pi}\int_L e^{-\ii\omega t}g_\omega\dd\omega.\end{equation}其中 $L$ 在全部奇点之上。证明终值定理 $\lim_{t\to\infty}g(t)=\lim_{\omega\to0}(-\ii\omega g_\omega)$；说明“稳定”与不衰减振荡的例外。
\end{exercise}
\begin{solution}
编者独立解答。若 $g(t)\to g_\infty$ 且剩余 $h(t)=g(t)-g_\infty$ 可积，$g_\omega=\ii g_\infty/\omega+h_\omega$（先取 $\operatorname{Im}\omega>0$），故 $-\ii\omega g_\omega\to g_\infty$。等价地变换 $\dot g$ 并积分分部。系统“稳定”在此须指除原点简单极点外，所有响应奇点位于导致时间衰减的 $\operatorname{Im}\omega<0$ 半平面，且无不衰减连续谱残留；单说增长率非正仍不足以保证极限存在。若有实轴 $\omega=\pm\Omega\ne0$ 极点，$g$ 可永远振荡而无通常终值；可先减去对应极点的 $e^{\mp\ii\Omega t}$ 项再对非振荡剩余用定理，或报告时间平均／Abel 极限并明确其定义。
\end{solution}
